\documentclass[11pt,a4paper]{article}

\usepackage[utf8]{inputenc}
\usepackage[T1]{fontenc}
\usepackage[portuguese]{babel}
\usepackage{lmodern}
\usepackage[a4paper,margin=2.2cm]{geometry}
\usepackage[table]{xcolor}
\usepackage{graphicx}
\usepackage{booktabs}
\usepackage{tabularx}
\usepackage{longtable}
\usepackage{enumitem}
\usepackage{needspace}
\usepackage{titlesec}
\usepackage{fancyhdr}
\usepackage{float}
\usepackage{caption}
\usepackage[htt]{hyphenat}
\usepackage{xurl}
\usepackage[hidelinks]{hyperref}

\definecolor{verde700}{HTML}{1B3A2F}
\definecolor{bronze}{HTML}{AA8D73}
\definecolor{tinta}{HTML}{1F2420}
\definecolor{tintasuave}{HTML}{5B5347}
\definecolor{marfimcl}{HTML}{F8F6F1}
\definecolor{critico}{HTML}{A4161A}
\definecolor{alto}{HTML}{D9480F}
\definecolor{medio}{HTML}{B08900}
\definecolor{baixo}{HTML}{3B6E8F}
\definecolor{ok}{HTML}{2B7A3D}

\newcommand{\sev}[2]{\colorbox{#1}{\color{white}\footnotesize\bfseries\strut\,#2\,}}
\newcommand{\critico}{\sev{critico}{CRÍTICO}}
\newcommand{\altoo}{\sev{alto}{ALTO}}
\newcommand{\medio}{\sev{medio}{MÉDIO}}
\newcommand{\baixo}{\sev{baixo}{BAIXO}}
\newcommand{\okk}{\sev{ok}{OK}}
\newcommand{\shot}[3][0.92]{%
  \begin{figure}[!htbp]\centering
  \fcolorbox{gray!40}{white}{\includegraphics[width=#1\linewidth]{screenshots/#2}}
  \caption{#3}\label{fig:#2}\end{figure}}
\newcommand{\bloco}[1]{\par\needspace{3\baselineskip}\textbf{\color{verde700}#1}\par}

\titleformat{\section}{\Large\bfseries\color{verde700}}{\thesection}{0.8em}{}
\titleformat{\subsection}{\large\bfseries\color{tinta}}{\thesubsection}{0.6em}{}
\captionsetup{font=small,labelfont=bf,labelsep=period}
\setlist{itemsep=2pt,topsep=4pt}
\setlength{\parskip}{4pt}
\setlength{\parindent}{0pt}
\setlength{\emergencystretch}{3em}
\pagestyle{fancy}
\fancyhf{}
\lhead{\small Carreto Lages --- domínio carretolages.com}
\rhead{\small Relatório de segurança e migração}
\cfoot{\small \thepage}

\begin{document}

\begin{titlepage}
  \vspace*{3cm}
  {\color{bronze}\small\scshape Relatório técnico --- confidencial}\\[0.6em]
  {\raggedright\Huge\bfseries\color{verde700} Domínio carretolages.com\\[0.2em] segurança e passagem\\ para a Cloudflare\par}\vspace{1.2em}
  {\large O que foi encontrado no site e no email atuais, os riscos, como resolver\\ e os passos para alojar o novo site na Cloudflare}\\[3em]
  \begin{tabular}{@{}ll}
    Data da análise: & 6 de novembro de 2026 \\
    Domínios: & \texttt{carretolages.com} (principal) e \texttt{carretolages.eu} \\
    Site atual: & \url{https://www.carretolages.com} (WordPress, alojado na PTisp) \\
    Site novo: & \url{https://carretolages.pages.dev} (pré-visualização) \\
  \end{tabular}
  \vfill
  {\small\color{tintasuave} Análise feita apenas com informação pública: pedidos normais a páginas do site,
  consultas DNS e serviços públicos de reputação. Não houve tentativas de entrar no site, no painel de
  alojamento nem no email. Este documento descreve falhas de segurança ainda abertas: deve ser partilhado
  só com quem vai tratar delas.}
\end{titlepage}

\tableofcontents
\clearpage

% =====================================================================
\section{Resumo}

\textbf{O site atual do escritório foi atacado e continua comprometido.} Desde outubro de 2023, alguém
publicou no site 5 artigos de spam: programas pirateados (``crack'' de Photoshop, Ableton, Adobe Premiere)
e réplicas de relógios Rolex. Esses artigos aparecem na secção ``Destaques'', com o cabeçalho, o logótipo e
os contactos do escritório. Um dos artigos foi publicado por uma conta de administrador chamada
\texttt{admin} que não pertence a ninguém do escritório nem à empresa que fez o site.

O site corre em software com 2 a 10 anos de atraso e está no \textbf{mesmo servidor do email}. Quem controla
o site pode, potencialmente, chegar ao alojamento e ao email.

\textbf{Boas notícias:} o Google não marca o site como perigoso, o servidor e o domínio não estão em listas
negras de spam e o email tem as proteções principais (DKIM e DMARC) bem configuradas.

\textbf{O que recomendamos:}
\begin{enumerate}
  \item \textbf{Esta semana:} mudar todas as palavras-passe (PTisp, cPanel, email, WordPress), apagar a conta
        \texttt{admin} e os artigos de spam e pedir à PTisp uma análise de malware ao alojamento
        (secção~\ref{sec:imediato}).
  \item \textbf{Substituir o site} pelo site novo, alojado na Cloudflare Pages: é estático, não tem base de
        dados, painel de login nem PHP, por isso este tipo de ataque deixa de ser possível.
  \item \textbf{Manter o email na PTisp} exatamente como está e passar apenas a gestão do DNS e o site para
        a Cloudflare (secção~\ref{sec:migracao}).
  \item \textbf{Desligar o WordPress} depois de o site novo estar no ar.
\end{enumerate}

\subsection*{Quadro de riscos}
\renewcommand{\arraystretch}{1.25}
\begin{longtable}{@{}p{0.75cm}p{7.15cm}p{2.2cm}p{4.6cm}@{}}
\toprule
\textbf{\#} & \textbf{Problema} & \textbf{Gravidade} & \textbf{Resolução} \\ \midrule
\endhead
\ref{r:spam} & Site comprometido: 5 artigos de spam publicados em nome do escritório & \critico & Apagar já; site novo \\
\ref{r:admin} & Conta de administrador desconhecida (\texttt{admin}) & \critico & Apagar; mudar todas as palavras-passe \\
\ref{r:software} & WordPress, PHP e plugins sem atualizações de segurança & \critico & Site novo; desligar WordPress \\
\ref{r:partilha} & Site e email no mesmo servidor e na mesma conta & \altoo & Palavras-passe, 2FA, rever email \\
\ref{r:exposicao} & Utilizadores, pastas e versões visíveis a qualquer pessoa; login aberto & \medio & Desaparece com o site novo \\
\ref{r:demo} & Conteúdo de demonstração do tema publicado como se fosse real & \medio & Desaparece com o site novo \\
\ref{r:links} & Ligações para domínios que já não existem (``Parceiros'') & \medio & Retirar; o site novo não as tem \\
\ref{r:cabecalhos} & Sem cabeçalhos de segurança no site & \baixo & O site novo já os tem \\
\ref{r:dns} & Pormenores de DNS e email (SPF, DNSSEC, \texttt{.eu}) & \baixo & Corrigir na migração \\
\ref{r:bom} & Reputação, certificado, DKIM/DMARC, registo do domínio & \okk & Manter \\
\bottomrule
\end{longtable}

% =====================================================================
\section{O que foi encontrado}
\label{sec:achados}

Cada ponto indica o que foi encontrado, porque é perigoso e como se resolve. As capturas de ecrã foram
feitas a 6 de novembro de 2026. As páginas com spam foram abertas com o JavaScript desligado, para não
executar código que lá possa estar.

% ---------------------------------------------------------------------
\subsection[Artigos de spam publicados no site]{Artigos de spam publicados no site \hfill\critico}
\label{r:spam}

\bloco{O que encontrámos}
O mapa do site (\texttt{wp-sitemap.xml}), que o WordPress entrega aos motores de busca, lista 9 artigos.
\textbf{5 são spam} colocado por terceiros:

\begin{center}\small
\begin{tabular}{@{}lll@{}}
\toprule
\textbf{Data} & \textbf{Endereço} & \textbf{Conta usada} \\ \midrule
22/10/2023 & \texttt{/ableton-10-1-6-mac-crack/} & n.º 1 \\
28/11/2024 & \texttt{/adobe-premiere-pro-transitions-pack-crack/} & n.º 101012 (\texttt{admin}) \\
11/12/2024 & \texttt{/replica-rolex-gmt-master-ii-126711-\ldots/} & n.º 1 \\
16/12/2024 & \texttt{/ableton-11-3-4-crack/} & n.º 1 \\
17/12/2024 & \texttt{/photoshop-download-with-crack/} & n.º 1 \\
\bottomrule
\end{tabular}
\end{center}

Os artigos usam o design do site, com o cabeçalho, o telefone e o email do escritório
(figuras~\ref{fig:03-spam-photoshop.png} e~\ref{fig:04-spam-rolex.png}).

\bloco{Porque é perigoso}
\begin{itemize}
  \item \textbf{Reputação:} um escritório de advogados a promover software pirateado e contrafação de
        marcas, com o seu nome e contactos.
  \item \textbf{Motores de busca:} este tipo de spam existe para usar a reputação do domínio e subir
        páginas de terceiros no Google. Se o Google o detetar, pode penalizar o domínio inteiro ou mostrar
        avisos nos resultados.
  \item \textbf{Sinal de controlo:} quem publicou estes artigos teve acesso de escrita ao site durante pelo
        menos 14 meses e pode ainda ter.
\end{itemize}

\bloco{Como resolver}
\begin{itemize}
  \item Apagar já os 5 artigos (e esvaziar o lixo do WordPress).
  \item Com o site novo, os endereços de spam passam a responder ``página não encontrada'' (404). É isso que
        leva o Google a esquecê-los: o mapa de redirecionamentos do site novo não os inclui de propósito.
  \item No Google Search Console, pedir a remoção destes endereços e consultar ``Problemas de segurança''
        (secção~\ref{sec:depois}).
\end{itemize}

\shot{02-sitemap-posts.png}{Mapa do site público: 5 dos 9 artigos são spam (cracks e réplicas).}
\shot{03-spam-photoshop.png}{Artigo de spam ``How to Get Photoshop for Free'' com o cabeçalho e os contactos do escritório.}
\shot{04-spam-rolex.png}{Artigo de spam a promover réplicas de relógios Rolex.}

% ---------------------------------------------------------------------
\subsection[Conta de administrador desconhecida]{Conta de administrador desconhecida \hfill\critico}
\label{r:admin}

\bloco{O que encontrámos}
A interface pública do WordPress (\texttt{/wp-json/wp/v2/users}) lista três contas
(figura~\ref{fig:06-api-users.png}):
\begin{itemize}
  \item n.º 1, uma conta pessoal;
  \item n.º 2, \texttt{TecniMobile Lda}, a empresa que fez o site;
  \item \textbf{n.º 101012, \texttt{admin}}.
\end{itemize}
O WordPress numera as contas por ordem (1, 2, 3\ldots). Um número como 101012 não resulta do uso normal:
é típico de contas inseridas diretamente na base de dados por um atacante, para manter o acesso mesmo que
as palavras-passe mudem. Foi esta conta que publicou o artigo de 28/11/2024
(figura~\ref{fig:05-api-posts.png}). Os outros 4 artigos de spam aparecem em nome da conta n.º 1, o que
indica que também essa conta, ou a base de dados, está nas mãos do atacante.

\bloco{Porque é perigoso}
Quem tem uma conta de administrador controla o site inteiro: pode publicar, instalar plugins com código
próprio, alterar formulários para recolher dados de visitantes e, através do código, chegar a ficheiros do
alojamento.

\bloco{Como resolver}
\begin{itemize}
  \item Apagar a conta \texttt{admin} (n.º 101012) e atribuir os conteúdos a outra conta.
  \item Mudar as palavras-passe de todas as contas do WordPress e da base de dados, e regenerar as chaves de
        segurança do \texttt{wp-config.php}. Isto termina as sessões abertas.
  \item Como a conta foi provavelmente criada na base de dados, assumir que o atacante pode ter deixado
        outros acessos (ficheiros PHP escondidos). Por isso a solução definitiva é \textbf{desligar o
        WordPress} e usar o site novo, que não tem base de dados nem login.
\end{itemize}

\shot[0.95]{06-api-users.png}{Lista pública de utilizadores: a conta \texttt{admin} com o número anómalo 101012.}
\shot[0.95]{05-api-posts.png}{Lista pública de artigos, com a data e a conta que publicou cada um.}

% ---------------------------------------------------------------------
\subsection[Software sem atualizações de segurança]{Software sem atualizações de segurança \hfill\critico}
\label{r:software}

\bloco{O que encontrámos}
\begin{center}\small
\begin{tabular}{@{}llll@{}}
\toprule
\textbf{Componente} & \textbf{Instalado} & \textbf{Atual} & \textbf{Estado} \\ \midrule
WordPress & 6.1.7 & 7.1.4 & marcado ``insecure'' pelo próprio WordPress \\
PHP (servidor) & 7.4.33 & 8.x & sem atualizações de segurança desde novembro de 2022 \\
WPBakery Page Builder & 4.11.1 & --- & versão de 2016 \\
Cookie Notice & 2.4.7 & 3.1.18 & desatualizado \\
jQuery & 3.6.1 & --- & desatualizado \\
Tema & ``lawyers'' & --- & tema comercial de demonstração \\
\bottomrule
\end{tabular}
\end{center}
A classificação ``insecure'' vem do serviço oficial do WordPress que diz a cada site se a sua versão é
segura (\url{https://api.wordpress.org/core/stable-check/1.0/}). O servidor também anuncia a versão do PHP
em cada resposta (\texttt{X-Powered-By: PHP/7.4.33}).

\bloco{Porque é perigoso}
Versões antigas têm falhas públicas e documentadas, com ferramentas automáticas que as procuram em milhares
de sites por dia. É muito provável que tenha sido assim que o site foi atacado.

\bloco{Como resolver}
Atualizar tudo seria possível, mas o tema e o WPBakery de 2016 dificilmente funcionam em versões atuais, e
a atualização não remove acessos que o atacante já tenha deixado. \textbf{Recomendamos não investir no
WordPress}: o site novo substitui-o e o WordPress é desligado logo a seguir à mudança. Até lá, o mínimo é
fazer o que está na secção~\ref{sec:imediato}.

\shot{09-readme.png}{Ficheiro \texttt{readme.html} do WordPress, acessível a qualquer pessoa.}

% ---------------------------------------------------------------------
\subsection[Site e email no mesmo servidor]{Site e email no mesmo servidor \hfill\altoo}
\label{r:partilha}

\bloco{O que encontrámos}
O site, o email (\texttt{mail}, \texttt{webmail}, \texttt{autodiscover}) e o painel de alojamento
(\texttt{cpanel}, \texttt{whm}, \texttt{ftp}) estão todos no mesmo servidor da PTisp, \texttt{130.185.84.194}.
O registo SPF do domínio, que diz que servidores podem enviar email em nome de
\texttt{@carretolages.com}, autoriza esse servidor.

\bloco{Porque é perigoso}
\begin{itemize}
  \item Código malicioso no site pode enviar email a partir desse servidor, e esse email pode passar as
        verificações anti-spam dos destinatários por vir de um servidor autorizado. É o cenário ideal para
        burlas com faturas ou pedidos de transferência ``em nome do escritório''.
  \item Se o atacante obteve credenciais do alojamento (cPanel), pode ler caixas de correio, criar contas de
        email ou criar reencaminhamentos silenciosos que copiam o correio para fora.
\end{itemize}

\bloco{Como resolver}
\begin{itemize}
  \item Mudar as palavras-passe da área de cliente PTisp, do cPanel, do FTP e de \textbf{todas} as caixas de
        email. Ativar a autenticação em dois passos (2FA) na área de cliente PTisp.
  \item No cPanel, verificar se há contas de email, reencaminhadores (``Forwarders'') ou filtros que ninguém
        do escritório criou. No webmail de cada pessoa, verificar regras de reencaminhamento.
  \item Pedir à PTisp uma análise de malware ao alojamento e os registos de acesso desde 2023.
  \item Com o site na Cloudflare e o WordPress desligado, o servidor da PTisp fica apenas com o email.
        Pode valer a pena perguntar à PTisp por um plano só de email.
\end{itemize}

% ---------------------------------------------------------------------
\subsection[Informação exposta e login aberto]{Informação exposta e login aberto \hfill\medio}
\label{r:exposicao}

\bloco{O que encontrámos}
\begin{itemize}
  \item \textbf{Nomes de utilizador públicos:} \texttt{/wp-json/wp/v2/users} e \texttt{/?author=1}
        revelam os nomes de login, que são metade do que é preciso para entrar.
  \item \textbf{Listagem de pastas:} \texttt{/wp-content/uploads/} e \texttt{/wp-includes/} mostram o
        conteúdo das pastas, como um explorador de ficheiros (figuras~\ref{fig:07-dir-uploads.png}
        e~\ref{fig:08-dir-wp-includes.png}).
  \item \textbf{Página de login aberta:} \texttt{/wp-login.php} aceita tentativas de qualquer lado, sem
        limite visível (figura~\ref{fig:10-wp-login.png}).
  \item \textbf{Versões anunciadas:} \texttt{readme.html}, a etiqueta \texttt{generator} e o cabeçalho
        \texttt{X-Powered-By} dizem exatamente que versões atacar.
\end{itemize}

\bloco{Porque é perigoso}
Juntas, estas informações permitem tentar palavras-passe de forma automática contra contas conhecidas e
escolher as falhas certas para as versões instaladas.

\bloco{Como resolver}
Tudo isto desaparece com o site novo, que é composto só por ficheiros estáticos, sem login, sem base de
dados e sem listagem de pastas. Até à mudança, a PTisp pode desligar a listagem de pastas (``Indexes'' no
cPanel) e limitar o acesso a \texttt{wp-login.php}.

\shot{07-dir-uploads.png}{Listagem pública da pasta de ficheiros carregados.}
\shot{08-dir-wp-includes.png}{Listagem pública da pasta interna \texttt{wp-includes}.}
\shot{10-wp-login.png}{Página de login do WordPress acessível a qualquer pessoa.}

% ---------------------------------------------------------------------
\subsection[Conteúdo de demonstração publicado]{Conteúdo de demonstração publicado \hfill\medio}
\label{r:demo}

\bloco{O que encontrámos}
Continuam publicadas páginas que vieram com o tema e nunca foram substituídas:
\begin{itemize}
  \item 4 artigos em inglês sobre ``personal injury lawyer'' e ``car accident'';
  \item 4 ``resultados de casos'' fictícios (``Auto Accident'', ``Medical Malpractice'', ``Domestic
        Violence'', ``Motocycle Accident'');
  \item 4 testemunhos de exemplo;
  \item páginas ``Página exemplo'', ``2 Columns'' e ``page2-2''.
\end{itemize}

\bloco{Porque é perigoso}
Estas páginas são indexadas como sendo do escritório e apresentam casos e testemunhos que não existem. Isso
prejudica a credibilidade e pode ser incompatível com as regras de publicidade da Ordem dos Advogados.

\bloco{Como resolver}
O site novo não tem estas páginas. Os endereços antigos com algum valor redirecionam para a página nova
equivalente (anexo~\ref{sec:redirects}). \textbf{Atenção:} o site novo também mostra, por decisão do
escritório, conteúdo de exemplo (testemunhos, casos e artigos) com o selo ``Exemplo ilustrativo''. Esse
conteúdo deve ser substituído por conteúdo real ou retirado logo que possível, pelas mesmas razões.

\shot{12-demo-post.png}{Artigo de demonstração do tema, em inglês, publicado em nome do escritório.}
\shot{11-demo-case.png}{``Resultado de caso'' de demonstração (``Auto Accident'').}

% ---------------------------------------------------------------------
\subsection[Ligações para domínios abandonados]{Ligações para domínios abandonados \hfill\medio}
\label{r:links}

\bloco{O que encontrámos}
A secção ``Parceiros'' da página inicial liga para \texttt{cristina-lanca-advogada.com} e
\texttt{jpalmslawfirm.com} (figura~\ref{fig:15-dead-partner-links.png}). Nenhum dos dois domínios está
registado (consulta ao registo \texttt{.com} a 6/11/2026).

\bloco{Porque é perigoso}
Qualquer pessoa pode registar estes domínios por poucos euros e pôr lá o que quiser, por exemplo uma página
falsa ou malware, e o site do escritório continua a enviar visitantes para lá.

\bloco{Como resolver}
Retirar as ligações. O site novo não as inclui.

\shot{15-dead-partner-links.png}{Secção ``Parceiros'': as ligações assinaladas apontam para domínios que já não existem (o endereço de cada ligação foi escrito a vermelho na captura, para se ler).}

% ---------------------------------------------------------------------
\subsection[Sem cabeçalhos de segurança]{Sem cabeçalhos de segurança \hfill\baixo}
\label{r:cabecalhos}

\bloco{O que encontrámos}
O site atual não envia nenhum dos cabeçalhos HTTP que os navegadores usam para limitar ataques.

\begin{center}\small
\begin{tabular}{@{}lcc@{}}
\toprule
\textbf{Cabeçalho} & \textbf{Site atual} & \textbf{Site novo} \\ \midrule
\texttt{Strict-Transport-Security} (HTTPS obrigatório) & não & sim \\
\texttt{Content-Security-Policy} (bloqueia scripts externos) & não & sim, estrita \\
\texttt{X-Frame-Options} / \texttt{frame-ancestors} (anti-clickjacking) & não & sim \\
\texttt{X-Content-Type-Options} & não & sim \\
\texttt{Referrer-Policy} & não & sim \\
\texttt{Permissions-Policy} (câmara, microfone, localização) & não & sim \\
Anuncia versão do PHP (\texttt{X-Powered-By}) & sim & não \\
\bottomrule
\end{tabular}
\end{center}

\bloco{Como resolver}
O site novo já tem todos estes cabeçalhos (ficheiro \texttt{public/\_headers}), testados automaticamente.

% ---------------------------------------------------------------------
\subsection[DNS e email: pormenores a corrigir]{DNS e email: pormenores a corrigir \hfill\baixo}
\label{r:dns}

\bloco{O que encontrámos}
\begin{itemize}
  \item \textbf{SPF de \texttt{carretolages.com}:}\\
        {\small\texttt{v=spf1 +a +mx +include:\_spf.cleanmx.pt ip4:130.185.84.194 -all}}\\
        O \texttt{+a} autoriza
        ``o servidor do site''. Depois da mudança, o site passa a ser a Cloudflare, e o \texttt{+a} passaria
        a autorizar os servidores da Cloudflare a enviar email em nome do escritório. Tem de sair. O IP do
        servidor de email já está escrito no registo.
  \item \textbf{SPF de \texttt{carretolages.eu}:} autoriza um endereço de rede privada
        (\texttt{192.168.1.90}), vários IPs externos sem explicação (um deles repetido) e termina em
        \texttt{\textasciitilde all} (``aceitar com desconfiança''). Deve ser limpo.
  \item \textbf{\texttt{mail.carretolages.com}} é um ``alias'' do domínio principal. Na mudança, se ficar
        assim, passa a apontar para a Cloudflare e o Outlook e os telemóveis deixam de receber email.
        \textbf{É o ponto mais delicado da migração.}
  \item \textbf{Certificado do email:} o servidor da PTisp usa um certificado
        \texttt{*.carretolages.com} válido até 3/01/2027. Os certificados com \texttt{*} são renovados
        através do DNS. Com o DNS na Cloudflare, é preciso confirmar com a PTisp que a renovação continua a
        funcionar para \texttt{mail} e \texttt{webmail}. Caso contrário, os programas de email mostram
        avisos de certificado a partir de janeiro.
  \item \textbf{DNSSEC} (assinatura do DNS contra falsificação) não está ativo.
\end{itemize}

\bloco{Como resolver}
Tudo é tratado nos passos da secção~\ref{sec:migracao}.

% ---------------------------------------------------------------------
\subsection[O que está bem]{O que está bem \hfill\okk}
\label{r:bom}

\begin{itemize}
  \item \textbf{Google Safe Browsing:} ``Nenhum conteúdo inseguro encontrado'' (figura~\ref{fig:13-safe-browsing.png}).
  \item \textbf{Listas negras:} nem o servidor nem o domínio constam das listas consultadas (Spamhaus ZEN e
        DBL, SpamCop, Barracuda, UCEPROTECT, PSBL, SURBL).
  \item \textbf{Email autenticado:} DKIM (seletor \texttt{default}) e DMARC com política
        \texttt{p=reject} e alinhamento estrito, a configuração mais protetora. Têm de ser copiados
        exatamente na migração.
  \item \textbf{Sem ``cloaking'':} o Google vê o mesmo conteúdo que um visitante (comparação das respostas
        para o robô do Google e para um navegador).
  \item \textbf{Certificado HTTPS} válido (Let's Encrypt, até 3/01/2027) para \texttt{.com} e \texttt{.eu};
        \texttt{http://} e o domínio sem \texttt{www} redirecionam para \texttt{https://www.}
  \item \textbf{Registo do domínio:} registado desde 2007 na NetEarth (revendedor da PTisp), com bloqueio de
        transferência ativo, \textbf{expira a 6/11/2027}. Convém confirmar que a renovação automática está
        ligada.
\end{itemize}

\shot{13-safe-browsing.png}{Google Safe Browsing: sem conteúdo inseguro detetado (6/11/2026).}

% =====================================================================
\section{O que fazer já, antes da mudança}
\label{sec:imediato}

Estes passos não dependem do site novo e devem ser feitos esta semana, pelo escritório ou pela PTisp.

\begin{enumerate}
  \item \textbf{Palavras-passe novas e diferentes} para a área de cliente PTisp (com 2FA), o cPanel, o FTP,
        todas as caixas de email e todas as contas do WordPress.
  \item \textbf{Apagar a conta \texttt{admin}} (n.º 101012) no WordPress.
  \item \textbf{Apagar os 5 artigos de spam} e esvaziar o lixo.
  \item \textbf{Verificar o email:} contas, reencaminhadores e filtros no cPanel; regras de
        reencaminhamento no webmail de cada pessoa.
  \item \textbf{Pedir à PTisp} uma análise de malware ao alojamento, a desativação da listagem de pastas e
        informação sobre acessos suspeitos desde 2023.
  \item \textbf{Não fazer} cópias de segurança do WordPress para reutilizar: estão contaminadas. Se for
        preciso arquivar textos ou imagens, copiar só esses conteúdos.
\end{enumerate}

% =====================================================================
\section{Passos para passar para a Cloudflare}
\label{sec:migracao}

\textbf{Princípio:} o domínio continua registado na PTisp e o email continua na PTisp. Passam para a
Cloudflare apenas o \textbf{DNS} (a ``lista telefónica'' do domínio) e o \textbf{site}. O plano gratuito da
Cloudflare chega. O custo mensal é 0~€.

\textbf{Porque o DNS e não só um registo:} a Cloudflare Pages só consegue servir o domínio sem \texttt{www}
quando gere o DNS. Também permite redirecionamentos, DNSSEC e regras sem depender do servidor antigo.

\subsection{Fase 0 --- Preparação (sem impacto no site nem no email)}
\begin{enumerate}
  \item \textbf{Conta Cloudflare do escritório}, criada com um email do escritório e com 2FA. O
        programador é convidado como membro, para o escritório nunca ficar dependente de uma conta pessoal.
  \item \textbf{Exportar a zona DNS completa} no cPanel (``Zone Editor'') ou pedi-la à PTisp em formato
        BIND (ficheiro de texto). A análise pública só vê os registos que se sabe procurar. A exportação
        garante que nada fica esquecido, em particular o DKIM, que é um texto longo e tem de ser copiado
        exatamente.
  \item \textbf{Publicar o site} na conta do escritório, com a versão para o domínio
        (indexável, com o conteúdo de exemplo):
\begin{verbatim}
npm run build:live
npx wrangler pages deploy dist --project-name carretolages --branch main
\end{verbatim}
        O endereço \texttt{carretolages.pages.dev} continua sempre com \texttt{noindex}, para não
        competir com o domínio no Google.
  \item \textbf{Combinar com a PTisp} a renovação do certificado de \texttt{mail} e \texttt{webmail}
        depois da mudança do DNS (secção~\ref{r:dns}).
\end{enumerate}

\subsection{Fase 1 --- Criar a zona na Cloudflare}
\begin{enumerate}
  \item Criar a zona, que fica ``pendente'' até à mudança dos nameservers:
\begin{verbatim}
cf zones create --name carretolages.com --type full
\end{verbatim}
        A resposta indica os \textbf{dois nameservers} atribuídos (do tipo \texttt{xxx.ns.cloudflare.com}).
  \item Importar a zona exportada na fase 0:
\begin{verbatim}
cf dns records import --zone carretolages.com --file carretolages.com.zone
\end{verbatim}
  \item \textbf{Rever registo a registo} e deixar a zona assim (``só DNS'' = nuvem cinzenta, sem passar pelo
        proxy da Cloudflare):
\end{enumerate}

\begin{center}\small
\renewcommand{\arraystretch}{1.2}
\begin{tabularx}{\linewidth}{@{}llXl@{}}
\toprule
\textbf{Nome} & \textbf{Tipo} & \textbf{Valor} & \textbf{Proxy} \\ \midrule
\texttt{www} & CNAME & \texttt{carretolages.pages.dev} (criado pela Pages) & sim \\
\texttt{@} & CNAME & \texttt{carretolages.pages.dev} (criado pela Pages) & sim \\
\texttt{@} & MX 10 & \texttt{mx1.cleanmx.pt} & --- \\
\texttt{@} & MX 20 & \texttt{mx2.cleanmx.pt} & --- \\
\texttt{@} & TXT & \texttt{v=spf1 +mx +include:\_spf.cleanmx.pt ip4:130.185.84.194 -all} & --- \\
\texttt{\_dmarc} & TXT & copiar exatamente (\texttt{v=DMARC1;p=reject;\ldots}) & --- \\
\texttt{default.\_domainkey} & TXT & copiar exatamente (DKIM) & --- \\
\textbf{\texttt{mail}} & \textbf{A} & \textbf{\texttt{130.185.84.194}} (hoje é CNAME: mudar!) & \textbf{não} \\
\texttt{webmail}, \texttt{autodiscover}, \texttt{autoconfig} & A & \texttt{130.185.84.194} & não \\
\texttt{cpanel}, \texttt{whm}, \texttt{webdisk}, \texttt{ftp} & A & \texttt{130.185.84.194} & não \\
\texttt{cpcalendars}, \texttt{cpcontacts} & A & \texttt{130.185.84.194} & não \\
outros da exportação (ex.: SRV) & --- & copiar como estão & não \\
\bottomrule
\end{tabularx}
\end{center}

Os registos antigos \texttt{@ A 130.185.84.194} e \texttt{www CNAME carretolages.com} saem, porque são
substituídos pelos do site novo. \textbf{Nada que tenha a ver com email pode ficar com proxy:} a Cloudflare
só encaminha tráfego web.

\subsection{Fase 2 --- Ligar o site ao domínio}
\begin{enumerate}
  \item Associar os dois endereços ao projeto Pages, que cria e valida os registos e os certificados:
\begin{verbatim}
cf pages projects domains create carretolages --name www.carretolages.com
cf pages projects domains create carretolages --name carretolages.com
\end{verbatim}
  \item Criar no painel da Cloudflare uma \textbf{regra de redirecionamento} (Rules $\rightarrow$ Redirect
        Rules, modelo ``Redirect from root to WWW''):\\
        \texttt{carretolages.com/*} $\rightarrow$ \texttt{https://www.carretolages.com/\$1}, código 301,
        mantendo o resto do endereço.
  \item Em SSL/TLS $\rightarrow$ Edge Certificates: ``Always Use HTTPS'' ligado e TLS mínimo 1.2.
\end{enumerate}
Até à fase 3, tudo isto fica preparado mas sem efeito, porque o mundo ainda consulta o DNS da PTisp.

\subsection{Fase 3 --- Mudar os nameservers (a mudança propriamente dita)}
\begin{enumerate}
  \item Na área de cliente da PTisp: Domínios $\rightarrow$ \texttt{carretolages.com} $\rightarrow$
        Nameservers. Substituir \texttt{nsp1}--\texttt{nsp4.mydnspt.net} pelos dois nameservers da
        Cloudflare.
  \item \textbf{Não apagar nada na PTisp.} Durante a propagação (até 24--48 horas), parte da internet
        consulta a PTisp e parte a Cloudflare. Como as duas têm os mesmos registos de email, o email
        funciona sempre. O site mostra o antigo ou o novo, conforme a pessoa.
  \item Confirmar que a zona ficou ativa:
\begin{verbatim}
cf zones get carretolages.com        # "status": "active"
\end{verbatim}
  \item \textbf{Testes no próprio dia:}
        \begin{itemize}
          \item enviar e receber um email de e para \texttt{geral@carretolages.com}, a partir de Gmail e
                de Outlook;
          \item abrir o email no telemóvel e no Outlook do escritório;
          \item abrir \url{https://www.carretolages.com} e \url{https://carretolages.com};
          \item abrir um endereço antigo, por exemplo \texttt{/contact/}, e confirmar que vai para
                \texttt{/contactos};
          \item abrir \texttt{/photoshop-download-with-crack/} e confirmar que dá ``página não
                encontrada''.
        \end{itemize}
\end{enumerate}

\textbf{Plano de recuo:} se algo falhar, repor os nameservers da PTisp na área de cliente. A zona antiga
continua intacta lá, por isso volta tudo ao estado anterior em poucas horas.

\subsection{Fase 4 --- Depois da mudança}
\label{sec:depois}
\begin{enumerate}
  \item \textbf{DNSSEC:} ativar na Cloudflare e colar o registo DS que a Cloudflare indica na área de
        cliente da PTisp:
\begin{verbatim}
cf dns dnssec edit --zone carretolages.com --status active
\end{verbatim}
  \item \textbf{Google Search Console:} verificar o domínio com um registo TXT, submeter
        \url{https://www.carretolages.com/sitemap-index.xml}, consultar ``Problemas de segurança'' e pedir a
        remoção dos 5 endereços de spam.
  \item \textbf{Desligar o WordPress} uma semana depois, se tudo estiver estável: apagar os ficheiros e a
        base de dados no cPanel. O plano PTisp fica só com o email.
  \item \textbf{Domínio \texttt{carretolages.eu}:} hoje redireciona através do servidor do WordPress. Para
        não deixar de funcionar quando o WordPress sair, criar também a zona \texttt{carretolages.eu} na
        Cloudflare com um redirecionamento para \url{https://www.carretolages.com}, mantendo os registos de
        email e limpando o SPF (secção~\ref{r:dns}).
  \item \textbf{Formulário de contacto:} o site novo ainda não envia mensagens. Para as entregar em
        \texttt{geral@carretolages.com} é preciso uma pequena função na Cloudflare e um serviço de envio
        de email configurado num subdomínio (por exemplo \texttt{envio.carretolages.com}), sem tocar no
        email principal. O Email Routing da Cloudflare não serve, porque obrigaria a passar o MX para a
        Cloudflare. A escolha do serviço de envio está pendente.
\end{enumerate}

% =====================================================================
\appendix
\section{Endereços antigos e para onde vão}
\label{sec:redirects}

O site novo inclui o ficheiro \texttt{public/\_redirects}, com redirecionamentos permanentes (301) dos
endereços do site antigo (com e sem a barra final). Assim, ligações em favoritos, emails antigos e
resultados do Google continuam a funcionar e passam a sua reputação para as páginas novas. Os 5 endereços
de spam ficam de fora de propósito.

\begin{center}\footnotesize
\renewcommand{\arraystretch}{1.1}
\begin{longtable}{@{}>{\raggedright\arraybackslash}p{8.4cm}>{\raggedright\arraybackslash}p{6.4cm}@{}}
\toprule
\textbf{Endereço antigo} & \textbf{Página nova} \\ \midrule
\endhead
\texttt{/about/} & \texttt{/sociedade} \\
\texttt{/historia/}, \texttt{/perfil-e-valores/}, \texttt{/pro-bono/} & \texttt{/sociedade/\ldots} (mesma página) \\
\texttt{/practice-areas/} & \texttt{/areas-de-atuacao} \\
\nolinkurl{/practice-items/direito-penal-e-das-contra-ordenacoes/} & \texttt{/areas-de-atuacao/direito-penal} \\
\texttt{/practice-items/direito-imobiliario-e-urbanismo/} & \texttt{/areas-de-atuacao/direito-imobiliario} \\
\texttt{/practice-items/direito-fiscal/} & \texttt{/areas-de-atuacao/direito-fiscal} \\
\texttt{/practice-items/direito-financeiro-bancario/} & \texttt{/areas-de-atuacao/direito-financeiro} \\
outras 5 áreas antigas (civil, administrativo, comercial, trabalho, seguros) & \texttt{/areas-de-atuacao} \\
\texttt{/lawyers/} & \texttt{/advogados} \\
\texttt{/team/jose-carreto-lages/}, \texttt{/team/joao-carlos-lages/}, \texttt{/team/catarina-pereira-carvalho/} & \texttt{/advogados/<mesmo nome>} \\
\texttt{/team/susana-portugal/}, \texttt{/team/rita-tavares/}, \texttt{/team/ana-mendes/} & \texttt{/advogados} \\
\texttt{/blog/} e os 4 artigos de demonstração & \texttt{/destaques} \\
\texttt{/case-results-2/} e \texttt{/cases/\ldots} & \texttt{/casos} \\
\texttt{/contact/} & \texttt{/contactos} \\
\texttt{/resolucao-de-litigios-online/} & \texttt{/resolucao-de-litigios} \\
\texttt{/faq/}, \texttt{/testimonials/}, \texttt{/testimonials-list/\ldots} & \texttt{/} (página inicial) \\
\texttt{/pagina-exemplo/}, \texttt{/2-columns/}, \texttt{/page2-2/} & \texttt{/} (página inicial) \\
\texttt{/cookies/}, \texttt{/termos-e-condicoes/}, \texttt{/politica-de-privacidade/} & mesma página (sem barra final) \\
5 artigos de spam & \textbf{404 --- não encontrado} \\
\bottomrule
\end{longtable}
\end{center}

\section{Como foi feita a análise}
\begin{itemize}
  \item \textbf{DNS:} consultas públicas (Cloudflare DNS-over-HTTPS e resolvedor 1.1.1.1) aos registos NS,
        A, MX, TXT, SPF, DKIM, DMARC, DS e subdomínios comuns de alojamento.
  \item \textbf{Registo do domínio:} RDAP do registo \texttt{.com} (Verisign).
  \item \textbf{Site:} pedidos normais (GET) às páginas públicas, ao mapa do site, à interface pública do
        WordPress e a ficheiros que o WordPress expõe por omissão. Não se tentou nenhum login e não se
        enviou nenhum formulário. Respostas comparadas entre um navegador e o robô do Google.
  \item \textbf{Versões:} serviço oficial do WordPress (\texttt{api.wordpress.org}) e diretório de plugins.
  \item \textbf{Reputação:} Google Safe Browsing e listas negras públicas por DNS.
  \item \textbf{Capturas:} Google Chrome automatizado (\texttt{scripts/domain-audit.mjs}), com JavaScript
        desligado nas páginas com spam. O servidor limita pedidos seguidos (erro 429), por isso as capturas
        foram feitas com pausas.
\end{itemize}

\end{document}
